\documentclass[10pt,a4paper]{amsart}
\usepackage[margin=19mm]{geometry}
\usepackage{amsmath,amssymb,mathtools}
\usepackage[scr=rsfs,cal=euler]{mathalfa}
\usepackage{xfrac}
\usepackage[unicode,hidelinks,bookmarks=false]{hyperref}
\newcommand{\ie}{i.e.\ }
\newcommand{\cf}{cf.\ }
\newcommand{\resp}{respectively\ }
\newcommand{\Subsection}{\subsection}
\providecommand{\nobreakdash}{\nobreak\hbox{-}}
\setlength{\emergencystretch}{2em}
\multlinegap0pt
\newtheorem{proposition}{Proposition}
\title{A 21-Coloring of the Plane Without Monochromatic Unit-Area Rectangles}
\author{Gary Hu, Yumo Liu}
\date{Source: arXiv:2609.00583v1 (CC BY 4.0)}
\begin{document}
\maketitle

\noindent\emph{Source and licence.} A 21-Coloring of the Plane Without Monochromatic Unit-Area Rectangles. Version arXiv:2609.00583v1 (CC BY 4.0), distributed by arXiv under the Creative Commons Attribution 4.0 International licence, http://creativecommons.org/licenses/by/4.0/, as recorded in the arXiv licence metadata for this exact version and shown on the abstract page https://arxiv.org/abs/2609.00583v1. Full licence notice: \texttt{licenses/arxiv-2609.00583-CC-BY-4.0.txt}.\par\smallskip

\noindent\textit{Editorial scope.} Counted unit: the article's proposition prop:sharp (source0.tex lines 102--104). The article's complete original proof of it is delivered with it (source0.tex lines 106--123, one \texttt{proof} environment of 18 lines). No mathematics is written, repaired or re-derived: the statement and the proof are the article's own lines, in the article's own notation, and no retained line is substituted or re-flowed.  No further source-local context is needed: every cross-reference of every retained line resolves inside the two delivered spans.  Nothing was omitted from any retained span, so the package carries no omission record.  No retained line contains a citation, so the wrapper carries no bibliography and no bibliography entry.  No retained line contains a figure environment, an image-inclusion command or drawing code, so no image is delivered and the package declares no asset dependency.  Editorial formatting only, in addition: an AMS \texttt{amsart} preamble that restates the article's own declarations and macros used by the retained lines, namely the environments \texttt{proposition} on separate counters, together with the standard packages those lines need.  The retained environments print in this document as Proposition 1 in order of appearance, whereas the article prints the same items with the numbering of its own full document.  The complete native source of the article as distributed in the arXiv e-print for this exact version is bundled unchanged as \texttt{sources/209041/source0.tex}.
\begin{proposition}\label{prop:sharp}
If $M\subset L$ has index at most $20$, then there exists $0\ne m\in M$ such that $d(m,4H_0)\le4/\sqrt3$.
\end{proposition}

\begin{proof}
In lattice coordinates, a point $x+y\omega$ lies in $4H_0$ if and only if $|2x+y|\le4$, $|x+2y|\le4$, and $|y-x|\le4$. Let $Q=4H_0-(1+\omega)/3$. A lattice point $p+q\omega$ lies in $Q$ if and only if $|2p+q+1|\le4$, $|p+2q+1|\le4$, and $|q-p|\le4$.
The lattice points in $Q$ are:

\begin{center}
\begin{tabular}{c|l}
$q$ & $p$ \\ \hline
$-3$ & $1$ \\
$-2$ & $-1,0,1,2$ \\
$-1$ & $-2,-1,0,1,2$ \\
$0$ & $-2,-1,0,1$ \\
$1$ & $-3,-2,-1,0,1$ \\
$2$ & $-2,-1$
\end{tabular}
\end{center}

There are $21$ lattice points in $Q$. If $[L:M]\le20$, two of them, say $x$ and $y$, lie in the same coset of $M$. Their difference $m=x-y$ is a nonzero point of $M$. Since $Q$ is a translate of $4H_0$, we have $m\in Q-Q=8H_0$. Hence $m/2\in4H_0$ and $|m|\le8/\sqrt3$. It follows that $d(m,4H_0)\le|m-m/2|=|m|/2\le4/\sqrt3$.
\end{proof}

\end{document}
